%==============================================================
\section{Πειραματικά Αποτελέσματα}
%==============================================================

Η παρούσα ενότητα παρουσιάζει την ποσοτική και ποιοτική αξιολόγηση του συστήματος AILS-NTUA. Αναλύεται η πειραματική διάταξη, παρατίθενται οι επίσημες επιδόσεις στα leaderboards του SemEval-2026 Task 8, και πραγματοποιείται εκτενής ανάλυση συνιστωσών (ablation studies) για τον προσδιορισμό της συνεισφοράς κάθε αρχιτεκτονικής επιλογής στα Tasks A, B και C.

%--------------------------------------------------------------
\subsection{Πειραματική Διάταξη}
%--------------------------------------------------------------

Η αξιολόγηση του συστήματος πραγματοποιήθηκε στο επίσημο benchmark του MTRAGEval (SemEval-2026 Task 8). Το σύνολο δεδομένων περιλαμβάνει $110$ πολύστροφες συνομιλίες που εκτείνονται σε συνολικά $842$ στροφές (turns) και καλύπτουν τέσσερα ετερογενή θεματικά πεδία (domains): \textit{FiQA} (οικονομικά), \textit{ClapNQ} (γενικές γνώσεις), \textit{Govt} (κυβερνητικά έγγραφα) και \textit{Cloud} (τεχνική υποστήριξη). Η ετερογένεια αυτή είναι σκόπιμη και δομικά σημαντική: όπως φαίνεται στον Πίνακα~\ref{tab:corpus-greek}, τα τέσσερα corpora διαφέρουν σημαντικά ως προς τη μέση πυκνότητα αποσπασμάτων ανά έγγραφο — από ατομικά έγγραφα-αναρτήσεις στο FiQA έως εκτενή εγκυκλοπαιδικά λήμματα στο ClapNQ — στρεσάροντας διαφορετικές πτυχές του μηχανισμού ανάκτησης. Η αξιολόγηση του Task A (Passage Retrieval) περιορίζεται στις $777$ στροφές που χαρακτηρίζονται ως απαντήσιμες (answerable) ή μερικώς απαντήσιμες (partially answerable).

\begin{table}[ht]
\centering
\small
\begin{tabular}{lcccc}
\toprule
\textbf{Χαρακτηριστικό} & \textbf{ClapNQ} & \textbf{FiQA} & \textbf{Govt} & \textbf{Cloud} \\
\midrule
Τύπος πεδίου & Εγκυκλοπαιδικό & Οικονομικό φόρουμ & Κυβερνητικά κείμενα & Τεχνική τεκμηρίωση \\
Σύνολο εγγράφων & 4{,}293 & 57{,}638 & 7{,}661 & 8{,}578 \\
Σύνολο αποσπασμάτων & 183{,}408 & 61{,}022 & 49{,}607 & 72{,}442 \\
Μέσος όρος αποσπ./έγγρ. & 42.7 & 1.1 & 6.5 & 8.4 \\
Κύρια πηγή & Wikipedia & StackExchange & Κυβ. ιστότοποι & IBM Cloud Docs \\
\bottomrule
\end{tabular}
\caption{Στατιστικά δομικής ετερογένειας του corpus MTRAG ανά θεματικό πεδίο.}
\label{tab:corpus-greek}
\end{table}

\paragraph{Υπερπαράμετροι και Ρύθμιση Συστήματος:}
Η βαθμονόμηση όλων των υπερπαραμέτρων πραγματοποιήθηκε αποκλειστικά στο development set. Στον αγωγό ανάκτησης, κάθε μία από τις $5$ LLM-based αναδιατυπώσεις εκδίδει ερώτημα προς τον ανακτητή \textbf{ELSER v1}, επιστρέφοντας τα κορυφαία $100$ αποσπάσματα κειμένου. Η επαναβαθμολόγηση εκτελείται μέσω του μοντέλου \textbf{Cohere Rerank v4} για τα $60$ επικρατέστερα έγγραφα, με σταθερά RRF $k=60$, ενώ η ιεραρχική συγχώνευση (Nested RRF) χρησιμοποιεί σταθερές εξομάλυνσης $k_{\text{internal}}=40$ στο Επίπεδο~1 και $k_{\text{final}}$ από $20$ έως $60$ ανά θεματικό πεδίο στο Επίπεδο~2.

\paragraph{Υπολογιστική Δρομολόγηση Μοντέλων (Model Routing):}
Για τη διαχείριση του οικονομικού κόστους των API, υιοθετήθηκε μια αρχιτεκτονική δρομολόγησης βάσει της κρισιμότητας του ρόλου (role-based model routing):
\begin{itemize}
    \item \textbf{DeepSeek-V3.2:} Χρησιμοποιείται για τις εργασίες υψηλής δομικής ακρίβειας, όπως η εξαγωγή των verbatim αποδεικτικών εκτάσεων (span extraction) και η υλοποίηση του Τεχνικού Κριτή (Technical Judge).
    \item \textbf{GPT-4o:} Αξιοποιείται για την Παραγωγή Υποψηφίων Απαντήσεων (Dual-Candidate Drafting) και για τους τρεις κριτές της πύλης απαντησιμότητας του Task C.
    \item \textbf{GPT-4o-mini:} Χρησιμοποιείται για υπολογιστικά ελαφριές εργασίες, όπως ο προκαταρκτικός έλεγχος απαντησιμότητας, ο Κριτής Ικανοποίησης Χρήστη και οι τελικές μικροδιορθώσεις.
\end{itemize}

%--------------------------------------------------------------
\subsection{Αποτελέσματα Ανάκτησης — Task A}
%--------------------------------------------------------------

\paragraph{Συγκριτική Ανάλυση Μοντέλων Ανάκτησης:}
Η αξιολόγηση εννέα ανακτητών (λεξιλογικών, πυκνών και learned sparse) στο development set, \textit{χωρίς} αναδιατύπωση ερωτήματος, κατέδειξε την ξεκάθαρη υπεροχή των learned sparse αναπαραστάσεων. Ο ανακτητής ELSER v1 πέτυχε $Recall@5: 0.483$ και $nDCG@5: 0.444$, υπερτερώντας του ισχυρότερου πυκνού ανακτητή (Cohere Embed v4, $R@5{=}0.397$) κατά $+8.6$ ποσοστιαίες μονάδες, και του κλασικού λεξιλογικού BM25 ($R@5{=}0.153$) κατά $+33.0$ ποσοστιαίες μονάδες στα ίδια $777$ ερωτήματα. Μέρος του πλεονεκτήματος αυτού μπορεί να οφείλεται στο ότι οι κρίσεις συνάφειας του MTRAG παρήχθησαν κυρίως από αποσπάσματα του ELSER. Το αποτέλεσμα είναι συνεπές με την υπόθεση ότι η επέκταση όρων (term expansion) σε επίπεδο corpus μειώνει το vocabulary gap χωρίς τον σημασιολογικό θόρυβο των cross-domain dense embeddings, αλλά δεν μπορεί να διαχωριστεί πλήρως από την παραπάνω μεροληψία των κρίσεων.

\paragraph{Ανάλυση Συνεισφοράς Αναδιατυπώσεων (Query Reformulation Ablation):}
Ο Πίνακας~\ref{tab:rewrite-ablation-greek} καταγράφει τη σωρευτική επίδραση της σταδιακής προσθήκης των πέντε στρατηγικών αναδιατύπωσης πάνω στον ανακτητή ELSER v1. Ο πλήρης συνδυασμός και των πέντε στρατηγικών υπό το αλγόριθμο Nested RRF αυξάνει το $Recall@5$ κατά \textbf{$+25.7\%$} σε σχέση με τη βασική γραμμή χωρίς αναδιατύπωση. Σε επίπεδο μεμονωμένης στρατηγικής, η \textit{Corpus-Specific} αναδεικνύεται ως η ισχυρότερη επιλογή στα πεδία ClapNQ ($R@5{=}0.637$) και FiQA ($R@5{=}0.491$), η \textit{Chain-of-Thought} υπερτερεί στο Govt ($R@5{=}0.580$), ενώ η \textit{Minimal} παραμένει η πλέον αποδοτική στο τεχνικό domain Cloud ($R@5{=}0.463$) — το HyDE, αντίθετα, αποδεικνύεται συστηματικά η ασθενέστερη μεμονωμένη στρατηγική σε όλα τα domains, χρήσιμη όμως ως συμπληρωματικό σήμα εντός του σχήματος Nested RRF. Ο πλήρης συνδυασμός ($R@5{=}0.607$) υπερβαίνει κατά $6.6$ μονάδες ακόμη και την καλύτερη μεμονωμένη στρατηγική· επειδή όμως περιλαμβάνει και την επαναβαθμολόγηση, η διαφορά είναι συνεπής με τη \textit{συμπληρωματικότητα} των στρατηγικών χωρίς να την απομονώνει.

\begin{table}[ht]
\centering
\small
\begin{tabular}{lcc}
\toprule
\textbf{Διαμόρφωση Ανάκτησης} & \textbf{Recall@5} & \textbf{Σχετική Βελτίωση ($\Delta$)} \\
\midrule
Μη ενισχυμένη βάση (αρχικό ερώτημα) & 0.483 & — \\
+ Ελάχιστη Αναδιατύπωση (Minimal) & 0.527 & $+9.1\%$ \\
+ Corpus-Ειδική Επέκταση & 0.558 & $+15.5\%$ \\
+ Αλυσίδα Σκέψης (CoT) & 0.573 & $+18.6\%$ \\
+ Υποθετικό Έγγραφο (HyDE) & 0.584 & $+20.9\%$ \\
+ Στόχευση Λέξεων-Κλειδιών & 0.591 & $+22.4\%$ \\
+ Σύντηξη Nested RRF (πλήρες σύστημα) & \textbf{0.607} & \textbf{$+25.7\%$} \\
\bottomrule
\end{tabular}
\caption{Σωρευτική ανάλυση συνεισφοράς (ablation) των πέντε στρατηγικών αναδιατύπωσης ερωτήματος στο development set. Οι γραμμές είναι σωρευτικές, οπότε κάθε αύξηση εξαρτάται από τη σειρά προσθήκης· ο πίνακας δεν απομονώνει τη συνεισφορά του reranker ή της ιεραρχικής σύντηξης.}
\label{tab:rewrite-ablation-greek}
\end{table}

\paragraph{Το Παράδοξο της Σύνθεσης Ανακτητών (Retriever Ensemble Paradox):}
Ένα από τα πιο σημαντικά ευρήματα της εργασίας είναι ότι η σύνθεση ετερογενών μοντέλων (π.χ. ELSER + BM25 + Dense μέσω επίπεδης RRF) υποβαθμίζει την ακρίβεια κορυφής (top-rank precision), παρότι βελτιώνει την οριακή ανάκληση σε μεγάλο βάθος. Όπως φαίνεται στον Πίνακα~\ref{tab:ensemble-paradox-greek}, στο πλήρως βελτιστοποιημένο στάδιο του αγωγού (ELSER με τις πέντε αναδιατυπώσεις), η προσθήκη όλων των εναλλακτικών ανακτητών αυξάνει το $Recall@100$ από $0.901$ σε $0.926$, ταυτόχρονα όμως \textit{μειώνει} το $Recall@5$ από $0.607$ σε $0.569$. Μια εξήγηση είναι ότι τα μοναδικά χρυσά έγγραφα που εισάγουν οι εναλλακτικοί ανακτητές εμφανίζονται σε μέση θέση από $27{,}5$ έως $54{,}2$ μετά τη σύντηξη ($39{,}5$ για όλους μαζί), κάτω δηλαδή από το όριο top-$10$, ενώ ταυτόχρονα η σύντηξη μετατοπίζει οριακά σχετικά αποτελέσματα του ELSER εκτός των κορυφαίων θέσεων. Η προτεινόμενη προσέγγιση \textit{ενιαίου ανακτητή με ποικιλότητα ερωτημάτων} αποφεύγει αυτή την παγίδα, καθώς όλα τα σήματα παραμένουν ευθυγραμμισμένα με τον ίδιο χώρο βαθμολόγησης. Σημειώνεται ότι οι κρίσεις συνάφειας του MTRAG παρήχθησαν κυρίως από αποσπάσματα που ανέκτησε το ELSER· αποσπάσματα που φέρνουν μόνο οι υπόλοιποι ανακτητές είναι πιθανότερο να μην έχουν κριθεί και να μετρώνται ως άσχετα. Η μεροληψία αυτή μπορεί από μόνη της να συμβάλλει στο παράδοξο, και δεν ποσοτικοποιήθηκε.

\begin{table}[ht]
\centering
\small
\begin{tabular}{lccc}
\toprule
\textbf{Διαμόρφωση (τελικό στάδιο, 5 αναδιατυπώσεις)} & \textbf{R@5} & \textbf{R@100} & \textbf{Μέση θέση μοναδικών χρυσών εγγρ.} \\
\midrule
ELSER v1 (μόνο) & \textbf{0.607} & 0.901 & — \\
+ Cohere Embed (dense) & 0.582 $\downarrow$ & 0.925 $\uparrow$ & 37.0 \\
+ SPLADE & 0.592 $\downarrow$ & 0.919 $\uparrow$ & 54.2 \\
+ BM25 & 0.536 $\downarrow$ & 0.915 $\uparrow$ & 27.5 \\
+ Όλοι οι ανακτητές & 0.569 $\downarrow$ & \textbf{0.926} $\uparrow$ & 39.5 \\
\bottomrule
\end{tabular}
\caption{Το παράδοξο σύνθεσης ανακτητών: η προσθήκη εναλλακτικών ανακτητών βελτιώνει το $R@100$ αλλά υποβαθμίζει συστηματικά το $R@5$.}
\label{tab:ensemble-paradox-greek}
\end{table}

\paragraph{Σταθερότητα ανά Στροφή Διαλόγου:}
Η απόδοση ανάκτησης εξαρτάται έντονα από τη θέση της στροφής (Σχήμα~\ref{fig:turn-depth-domain}). Χωρίς αναδιατύπωση, το $Recall@5$ πέφτει από $0.879$ στην πρώτη στροφή σε $0.428$ στις στροφές $6+$ (σχετική πτώση $51\%$)· με το πλήρες σύστημα η πτώση περιορίζεται στο $39\%$ (από $0.890$ σε $0.540$). Η πτώση δεν είναι σταδιακή: συγκεντρώνεται κυρίως μεταξύ της πρώτης στροφής, που είναι εξ ορισμού αυτοτελής, και της δεύτερης, ενώ στη συνέχεια οι καμπύλες μένουν σχεδόν σταθερές. Συνεπώς η επίδραση της θέσης οφείλεται κυρίως στο αν το ερώτημα είναι αυτοτελές ή εξαρτάται από προηγούμενες στροφές (αναφορές αντωνυμίας στο $97\%$ των μη-πρώτων στροφών, συγκριτικές ερωτήσεις-συνέχειες στο $67\%$). Η αναδιατύπωση περιορίζει, χωρίς να εξαλείφει, το χάσμα αυτό.

%--------------------------------------------------------------
\subsection{Αποτελέσματα Τεκμηριωμένης Παραγωγής — Task B}
%--------------------------------------------------------------

\paragraph{Ablation Analysis του Agentic Pipeline:}
Ο Πίνακας~\ref{tab:gen-ablation-greek} καταγράφει τη σωρευτική συνεισφορά κάθε σταδίου του πολυσταδιακού αγωγού παραγωγής στο Task B. Η σταδιακή ενσωμάτωση των συνιστωσών οδηγεί σε συνολική αύξηση του Αρμονικού Μέσου (HM) κατά \textbf{$+8.9$ μονάδες} ($0.654 \rightarrow 0.743$) έναντι μίας απλής, μονοβηματικής (greedy) παραγωγής.

\begin{table}[ht]
\centering
\small
\begin{tabular}{lc}
\toprule
\textbf{Διαμόρφωση Αγωγού} & \textbf{Αρμονικός Μέσος (HM)} \\
\midrule
Βάση: Μονοβηματική Greedy Παραγωγή ($T{=}0.0$) & 0.654 \\
+ Εξαγωγή Αποδεικτικών Εκτάσεων (Verbatim) & 0.688 \\
+ Παραγωγή Δύο Υποψηφίων (τυχαία επιλογή) & 0.706 \\
+ Τεχνικός Κριτής & 0.725 \\
+ Κριτής Ικανοποίησης Χρήστη ($60\%$ δειγματοληψία) & 0.738 \\
+ Τελικές Μικροδιορθώσεις & \textbf{0.743} \\
\bottomrule
\end{tabular}
\caption{Ανάλυση συνεισφοράς (ablation) του αγωγού παραγωγής Task B στο development set.}
\label{tab:gen-ablation-greek}
\end{table}

\paragraph{Η Κρισιμότητα του Evidence Span Extraction:}
Η εισαγωγή του σταδίου εξαγωγής αποδεικτικών εκτάσεων αναδεικνύεται ως η \textit{μεμονωμένα} μεγαλύτερη βελτίωση του αγωγού ($+0.034$ HM, Πίνακας~\ref{tab:gen-ablation-greek}), μαζί με αύξηση κατά $+0.041$ στη μετρική πιστότητας χωρίς αναφορά $RL_F$. Η verbatim απομόνωση των σχετικών προτάσεων περιορίζει την τάση των μοντέλων να ενσωματώνουν άσχετες πληροφορίες (distractors) από το υπόλοιπο σώμα των εγγράφων, εμποδίζοντας έτσι τον generator να «δραπετεύσει» από τα παρεχόμενα στοιχεία.

\paragraph{Αποτίμηση του Σχήματος Πολλαπλών Κριτών:}
Η εισαγωγή των δύο LLM κριτών (Technical, User Satisfaction) σε συνδυασμό με τον όρο Διαμόρφωσης Εκχύλισης (Extractiveness Shaping) στοχεύει στην εξισορρόπηση της πραγματολογικής ακρίβειας με τη συνομιλιακή φυσικότητα, και κάθε στάδιο προσθέτει θετική συνεισφορά στο ablation του Πίνακα~\ref{tab:gen-ablation-greek}. Ο όρος εκχύλισης ποινικοποιεί προσχέδια που είτε απομακρύνονται υπερβολικά από τα εξαχθέντα spans (κίνδυνος ψευδαίσθησης) είτε τα αντιγράφουν σχεδόν αυτολεξεί, ευνοώντας επικάλυψη 4-γραμμάτων εντός της ζώνης $[0.28, 0.38]$, που επιλέχθηκε στο development set (Πίνακας~\ref{tab:extractiveness-band-en}).

\paragraph{Οικονομική Αποδοτικότητα μέσω Routing:}
Η δρομολόγηση μοντέλων (DeepSeek-V3.2 για εξαγωγή/τεχνική κρίση, GPT-4o για τη σύνθεση, GPT-4o-mini για ελαφριές εργασίες) είναι σχεδιαστική επιλογή με κίνητρο το κόστος. Με ονομαστικές τιμές API, ο αγωγός του Task B κοστίζει κατά εκτίμηση $\$0{,}010$--$0{,}012$ ανά στροφή, από τα οποία τα $\$0{,}009$ αντιστοιχούν στις δύο κλήσεις του GPT-4o, έναντι εκτιμώμενων $\$0{,}024$ αν όλα τα στάδια χρησιμοποιούσαν GPT-4o, δηλαδή περίπου $2$--$2{,}4$ φορές χαμηλότερο κόστος.\footnote{Το δημοσιευμένο άρθρο~\citep{athanasiou2026ailsntua} αναφέρει $\$0{,}008$ ανά στροφή και $3\times$ μείωση. Το άθροισμα των εκτιμήσεων ανά στάδιο, με προσμέτρηση και των δύο κλήσεων παραγωγής, δίνει $\$0{,}010$--$0{,}012$, τιμή που υιοθετείται εδώ.} Δεν διατηρήθηκε ελεγχόμενη σύγκριση ποιότητας με ομοιόμορφη χρήση ενός μοντέλου, οπότε δεν διατυπώνεται ισχυρισμός για την επίδραση της δρομολόγησης στην ποιότητα.

%--------------------------------------------------------------
\subsection{Η Απαντησιμότητα ως Καθοριστικός Παράγοντας --- Task C}
%--------------------------------------------------------------

Η end-to-end αξιολόγηση (Task C) παρουσιάζει απότομη πτώση σε σχέση με την παραγωγή με αποσπάσματα αναφοράς ($HM$: $0.7698 \rightarrow 0.5409$ στο test set). Σε συμφωνία με το δημοσιευμένο άρθρο~\citep{athanasiou2026ailsntua}, θεωρούμε τη βαθμονόμηση της απαντησιμότητας καθοριστικό παράγοντα υπό την κατανομή του test set: η πύλη αποδέχεται τις περισσότερες μη απαντήσιμες στροφές ακόμη και στο development set και βαθμονομήθηκε σε ποσοστό μη απαντήσιμων $6.5\%$, ενώ στο test set αντιμετωπίζει $19.1\%$. Η ανάκτηση παραμένει πάντως σημαντικός παράγοντας. Στο Task C ο generator λαμβάνει μόνο τα τρία πρώτα αποσπάσματα, οπότε το αν τα αποδεικτικά στοιχεία φτάνουν σε αυτόν καθορίζεται από την ακρίβεια κορυφής ($Recall@5{=}0.607$ στο development set, Πίνακας~\ref{tab:ensemble-paradox-greek}) και όχι από την ανάκληση σε βάθος ($Recall@100{=}0.901$). Επιπλέον, τα σφάλματα ανάκτησης στις κορυφαίες θέσεις αντιστοιχούν στις περισσότερες αποτυχίες του development set που εξετάστηκαν χειροκίνητα (Πίνακας~\ref{tab:error-decomp-greek}), σε ένα σύνολο όπου οι μη απαντήσιμες στροφές είναι σπάνιες. Η ανάλυση των επίσημων σκορ του test set ανά κατηγορία απαντησιμότητας θα ποσοτικοποιούσε τις δύο συνεισφορές. Όπως φαίνεται στον Πίνακα~\ref{tab:answerability-greek}, καμία διαμόρφωση πύλης δεν ξεπερνά το $25\%$ $F1$ στην ανίχνευση μη απαντήσιμων στροφών στο πλήρες development set. Το προτεινόμενο σχήμα πολλαπλών κριτών πετυχαίνει $95.8\%$ ανάκληση στις απαντήσιμες στροφές αλλά μόλις $21.8\%$ στις μη απαντήσιμες, δηλαδή $12$ από τις $55$ (95\% διάστημα Wilson $12{,}9$--$34{,}4\%$), ένδειξη ισχυρής μεροληψίας αποδοχής (acceptance bias). Η ποιοτική εξέταση των αποφάσεων της πύλης δείχνει ότι οι κριτές Εγγράφου και Αποσπάσματος συχνά επισημαίνουν ανεπαρκή στοιχεία, ενώ ο Κριτής Απάντησης τείνει να αποδέχεται προσχέδια που παραμένουν εύλογα, δίνοντας προτεραιότητα στην ευλογοφάνεια έναντι της επάρκειας των αποδείξεων.

\begin{table}[ht]
\centering
\small
\begin{tabular}{lccc}
\toprule
\textbf{Μέθοδος} & \textbf{Ανάκληση Απαντ.} & \textbf{Ανάκληση Μη Απαντ.} & \textbf{F1 Μη Απαντ.} \\
\midrule
Πάντα Απάντηση (χωρίς πύλη) & $100.0\%$ (777/777) & $0.0\%$ (0/55) & $0.0\%$ \\
Αγωγός + Κανόνας Επαλήθευσης & $92.3\%$ (717/777) & $27.3\%$ (15/55) & $23.1\%$ \\
\textbf{Πολλαπλοί Κριτές (δικό μας)} & $95.8\%$ (744/777) & $21.8\%$ (12/55) & $24.0\%$ \\
\bottomrule
\end{tabular}
\caption{Μετρικές ανίχνευσης απαντησιμότητας στο development set ($832$ στροφές: $777$ απαντήσιμες ή μερικώς απαντήσιμες και $55$ μη απαντήσιμες· εξαιρούνται οι $10$ συνομιλιακές στροφές). Καμία διαμόρφωση που αξιολογήθηκε δεν ξεπερνά το $25\%$ $F1$ στις μη απαντήσιμες στροφές του πλήρους development set.}
\label{tab:answerability-greek}
\end{table}

\paragraph{Ποσοτική Κατανομή Σφαλμάτων και Εξαρτήσεις μεταξύ Στροφών:}
Ο συγγραφέας εξέτασε χειροκίνητα $50$ περιπτώσεις αποτυχίας του development set και απέδωσε σε καθεμία μία κύρια αιτία (Πίνακας~\ref{tab:error-decomp-greek}). Οι περιπτώσεις δεν επιλέχθηκαν με τεκμηριωμένη τυχαία διαδικασία και η κατηγοριοποίηση έγινε από έναν σχολιαστή, οπότε τα ποσοστά είναι ενδεικτικά.

\begin{table}[ht]
\centering
\small
\begin{tabular}{lcc}
\toprule
\textbf{Αιτία Σφάλματος} & \textbf{Ποσοστό} & \textbf{Κύριο Πεδίο} \\
\midrule
Σφάλμα ανάκτησης (περιεχόμενο τρέχουσας στροφής) & $38\%$ & FiQA \\
Σφάλμα ανάκτησης (ανεπίλυτη εξάρτηση από προηγούμενες στροφές) & $34\%$ & Cloud \\
Εσφαλμένη εκτίμηση Απαντησιμότητας & $18\%$ & Όλα \\
Καθαρό σφάλμα παραγωγής (σωστά αποσπάσματα) & $10\%$ & ClapNQ \\
\bottomrule
\end{tabular}
\caption{Κύρια αιτία σε $50$ περιπτώσεις αποτυχίας του Task C που εξετάστηκαν χειροκίνητα (development set, ένας σχολιαστής). Οι αιτίες που σχετίζονται με την ανάκτηση αντιστοιχούν στο $72\%$.}
\label{tab:error-decomp-greek}
\end{table}

Οι αιτίες που σχετίζονται με την ανάκτηση αντιστοιχούν συνολικά στο $72\%$ των περιπτώσεων και η εσφαλμένη εκτίμηση απαντησιμότητας στο $18\%$. Η δεύτερη κατηγορία ανάκτησης ($34\%$) χρειάζεται προσεκτική ανάγνωση: επειδή το MTRAGEval παρέχει το ιστορικό αναφοράς και όχι τις προηγούμενες απαντήσεις του συστήματος, ένα σφάλμα του συστήματος στη στροφή $t-1$ δεν μπορεί να περάσει στο $H_t$. Η κατηγορία $34\%$ αφορά επομένως στροφές των οποίων το ερώτημα εξαρτάται από προηγούμενες στροφές και η αναδιατύπωση δεν επέλυσε την εξάρτηση, με αποτέλεσμα να αποτύχει η ανάκτηση στη στροφή $t$ παρότι το ιστορικό είναι σωστό. Οι περιπτώσεις αυτές υποδεικνύουν την επίλυση συναναφοράς και έλλειψης ως σημαντική εναπομένουσα αδυναμία της ανάκτησης, σε συμφωνία με το χάσμα μεταξύ αυτοτελών και μη αυτοτελών ερωτημάτων.

%--------------------------------------------------------------
\subsection{Επίσημα Αποτελέσματα Leaderboard}
%--------------------------------------------------------------

Στην επίσημη διεθνή αξιολόγηση του SemEval-2026 Task 8, η αρχιτεκτονική \textbf{AILS-NTUA} κατέλαβε την 1η θέση στο Task A και τη 2η στο Task B (Πίνακας~\ref{tab:leaderboard-greek}):
\begin{itemize}
    \item \textbf{1η Θέση στο Task A (Passage Retrieval):} Το σύστημα πέτυχε την κορυφαία επίδοση με \textbf{$nDCG@5: 0.5776$}, σημειώνοντας διαφορά \textbf{$+20.5\%$} έναντι της ισχυρότερης επίσημης γραμμής βάσης (baseline) του διαγωνισμού, μεταξύ $38$ συμμετοχών.
    \item \textbf{2η Θέση στο Task B (Grounded Generation):} Το σύστημα κατέλαβε τη δεύτερη θέση με Αρμονικό Μέσο \textbf{$HM: 0.7698$} ($RB_{alg}{=}0.633$, $RL_F{=}0.897$, $RB_{llm}{=}0.832$), σε συμφωνία με τα ευρήματα για το span extraction και το multi-judge framework, μεταξύ $26$ συμμετοχών.
    \item \textbf{11η Θέση στο Task C (End-to-End RAG):} Το σύστημα πέτυχε $HM: 0.5409$ μεταξύ $29$ συμμετοχών, μόλις οριακά πάνω από την ισχυρότερη γραμμή βάσης των διοργανωτών ($0.5366$, $+0.004$). Η πτώση από το $HM{=}0.7698$ του Task B οφείλεται κυρίως στη μετρική $RB_{alg}$ ($0.633 \rightarrow 0.400$) και συμφωνεί με την ανάλυση της Ενότητας 1.5.4: η πύλη βαθμονομήθηκε σε development set με ποσοστό μη απαντήσιμων στροφών περίπου ίσο με το ένα τρίτο εκείνου του test set ($6.5\%$ έναντι $19.1\%$, Πίνακας~\ref{tab:dev-test-greek}), οπότε το κατώφλι $\tau{=}0.7$ δεν ταιριάζει στις συνθήκες του test set. Συμβάλλουν επίσης τα σφάλματα ανάκτησης στις κορυφαίες θέσεις, αφού στον generator φτάνουν μόνο τρία αποσπάσματα. Το αν ένα χαμηλότερο κατώφλι (π.χ. $\tau \approx 0.6$) θα απέδιδε καλύτερα στο test set αποτελεί υπόθεση που δεν μπορεί να ελεγχθεί με μία μόνο υποβολή· στο development set, το $\tau{=}0.6$ αποδίδει ελαφρώς χειρότερα από το $\tau{=}0.7$.
\end{itemize}

\begin{table}[ht]
\centering
\small
\begin{tabular}{llcc}
\toprule
\textbf{Task} & \textbf{Μετρική} & \textbf{Επίδοση} & \textbf{Κατάταξη} \\
\midrule
A — Ανάκτηση & nDCG@5 & 0.5776 & 1/38 \\
B — Παραγωγή & HM & 0.7698 & 2/26 \\
C — End-to-End RAG & HM & 0.5409 & 11/29 \\
\bottomrule
\end{tabular}
\caption{Επίσημα αποτελέσματα του AILS-NTUA στο test set του SemEval-2026 Task 8.}
\label{tab:leaderboard-greek}
\end{table}